\section{Introduction}
\label{sec:introduction}

\emph{Model distillation} transfers capabilities from a teacher model to a student model through teacher-generated supervised training data~\cite{hinton2015distilling}. In language models and autonomous agents, fine-tuning on teacher reasoning and tool-use trajectories allows the students to acquire complex behaviors at low cost~\cite{hsieh2023distilling, mukherjee2023orca, liu2026structured, kang2026distilling}. By systematically collecting outputs from proprietary closed-source models, competitors can acquire state-of-the-art capabilities while bypassing substantial research and training investments~\cite{anthropicDetectingPreventing}. Recent disclosures from leading AI providers indicate that unauthorized distillation campaigns now operate at an industrial scale~\cite{OpenAIkey, googleGTIGThreat, anthropicDetectingPreventing}, presenting serious intellectual property risks across the AI ecosystem.


% [narrow background: we should talk about some anti-distillation approaches documented by mainstream providers, and how they still get evaded. also,
% the providers report shows a key missing piece: although they can detect and stop unauthorized distillation campaigns through traffic monitoring, knowing which traces are potentially collected, they don't have a way to prove a suspect student model was trained on these output -- i..e, a forensic tracing technique is missing.
% ]
Mainstream providers (e.g., OpenAI, Anthropic and Google) combine \emph{traffic screening}, \emph{account controls}, and \emph{output safeguards} to counter unauthorized distillation~\cite{}. \mingxun{TODO: cite}
Traffic screening uses classifiers and manual investigation to identify suspicious collection.
Account controls verify users and block accounts associated with abuse.
Output safeguards restrict access to reasoning traces and other useful training material.
Distillers evade these controls through proxy networks that replace banned accounts and mix collection with ordinary traffic, as well as third-party routers that conceal the source of requests~\cite{anthropicDetectingPreventing,OpenAIkey}.
They may also clean the collected outputs before training~\cite{OpenAIkey}.
Once outputs have been collected, blocking further access does not prevent their subsequent use for training.

% Writing guidance: first describe fully passive comparisons between the suspect
% student and teacher; then describe active embedding of learnable signals.
These defenses help prevent collection and disrupt ongoing campaigns, but leave a gap in \emph{post-hoc tracing}.
Even when a provider identifies potentially collected trajectories, traffic evidence alone does not establish whether a suspect student was trained on them.
Existing approaches to post-hoc detection follow two main routes.
\emph{Fully passive methods} compare a suspect student model's behavior against the teacher under selected conditions, looking for similarities that may indicate distillation~\cite{yang2026agents,rawat2026reference,panfilov2026stealing}.
\emph{Active watermarking} methods instead embed learnable signals in teacher outputs and later probe the student to test whether it has inherited them~\cite{xu2026antidistillation,wang2026protecting,an2026sequential}.

% Writing guidance: acknowledge the appeal of passive and active methods, then
% explain the joint deployment requirements and state the research question.
% Provider reports do not establish deployment of these forensic methods at scale;
% do not infer that deployment never occurs or attribute non-adoption to a proven cause.
Both approaches offer a plausible basis for forensic tracing, but there is no clear evidence showing that they have been deployed in practice (noticably missing from the mainstream providers' report on anti-distillation efforts~\cite{}).
Passive comparisons must distinguish evidence of training use from similarities that unrelated models may share.
Our independent analysis and experiments identify concrete ambiguities in reasoning-prefill detection and AgentEcho (see Section~\ref{sec:intro-contributions}).
Active watermarks must survive distillation and data cleaning while preserving the teacher's task utility.
The research gap is to meet these requirements together: reliable detection must provide strong evidence tied to the provider's outputs, while keeping serving and auditing costs low.
The provider must obtain this evidence under \emph{black-box student access}, without knowing which collected trajectories entered training.
We therefore ask:

\begin{itemize}
    \item[] \emph{How can a provider obtain reliable evidence of distillation through black-box student queries when the training data are unknown and may have been cleaned?}
\end{itemize}

\subsection{Our Contributions}
\label{sec:intro-contributions}

Our contributions are twofold.

\partitle{Systematization and new empirical evidence}
Our first contribution is a systematic evaluation of existing anti-distillation techniques across their access requirements, operational costs, evidential strength, and resilience to data processing (Table~\ref{tab:anti-distillation-comparison}). We complement this taxonomy with new analytical and empirical evidence exposing fundamental ambiguities in prominent detection methods. For reasoning-prefill detection~\cite{panfilov2026stealing}, controlled comparisons on public problems show that shared conventional derivations (such as standard differential-equation integrating factors) and alternative prompt prefixes can produce substantial answer overlap without teacher distillation. For AgentEcho~\cite{yang2026agents}, we construct an analytical counterexample showing that an 11-node graph with a single edge and an 11-node complete 10-edge chain yield proportional feature vectors and an identical maximal dependency similarity ($S_{\mathtt{dep}} = 1.0$). Furthermore, our trajectory evaluations show that reasoning distillation concurrently increases behavioral similarity toward multiple unrelated models. These findings demonstrate that unwatermarked behavioral similarity cannot uniquely identify a claimed training source.

\partitle{Detection through retained behavioral evidence}
Our second contribution connects interventions during model serving with post-hoc verification of a suspect student. In our setting, the provider creates and privately records behavioral evidence while serving trajectories. Service-side traffic screening subsequently identifies an anomalous candidate set of suspicious interactions $\mathcal{S}$ for auditing. Crucially, the provider does not need to know which specific interactions entered the distiller's training set $\mathcal{D}$, nor assume that $\mathcal{D} \subseteq \mathcal{S}$. Instead, our main insight is to audit the suspect student using targeted behavioral probes derived from the retained evidence, relying on the statistical signal surviving in the non-empty intersection $\mathcal{S} \cap \mathcal{D}$ (Figure~\ref{fig:threat-model}).

\partitle{Embedding learnable auxiliary behaviors}
We instantiate this approach in \name{}, a behavioral watermarking method for language model agents (Figure~\ref{fig:overview}). Rather than altering text tokens or modifying core task actions, \name{} inserts safe \emph{auxiliary tool calls} into teacher trajectories immediately following normal core actions. Auxiliary calls are restricted to a pre-classified safe tool subset $\mathcal{F}_{\mathtt{safe}}$ and strictly bound to parameter values drawn from the interaction context, ensuring that they execute safely and appear natural. A keyed pseudo-random schedule and candidate scoring policy favor repeatable core--auxiliary action pairs, instilling consistent behavioral habits that the student learns during distillation. Sparse insertion limits additional interaction steps while preserving the teacher's core task execution and accuracy. Each successful insertion creates a \emph{private evidence card} recording the execution context, the core action, and the injected auxiliary parameters.

\partitle{Paired-probe verification and trace attribution}
To audit a suspect student, observing an auxiliary call in isolation is insufficient because an untrained model might occasionally invoke the tool naturally. \name{} resolves this ambiguity through \emph{paired probes} constructed from the retained evidence cards. For each card, the provider creates a \emph{real probe} preserving the recorded core action and a \emph{counterfactual control probe} that replaces the core action with a semantically equivalent alternative achieving the same goal. The provider queries the black-box student and tests whether it preferentially reproduces the target auxiliary behavior in real contexts over control contexts. Aggregating probe variants at the card level eliminates intra-card dependence, and the detection rule requires both a one-sided exact sign test ($p < 0.05$) and a minimum full-hit margin ($\ge 0.05$). The same evidence cards further support trace-level attribution, ranking candidate trajectories to identify specific samples included in the adversary's training set.

\partitle{Empirical evaluation}
We evaluate \name{} across three agent benchmarks (BFCL, SWE-bench, and Telecom) and eight teacher--student pairs spanning four student model families. In an audit setting where the provider inspects 100 suspicious trajectories of which only 50 were used for distillation, \name{} detects all 24 distilled student models ($p < 10^{-3}$ to $10^{-14}$) while producing zero false positives across 48 clean-model evaluations (all $p \ge 0.062$). Watermark embedding preserves teacher utility: task accuracy changes by at most 3.0 percentage points, with mean interaction length increasing by only 0.60 steps. Furthermore, \name{} remains effective under aggressive data processing, successfully detecting students after 1:10 dilution with clean data, mixing across three independent teachers, semantic rewriting, trajectory truncation, and adaptive LLM sanitization explicitly targeting auxiliary behaviors.
